%%%%%%%%%%%%%%%%%%%%%%%%%%%%%%%%%%%%%%%%%%%%%%%%%%%%%%%%%%%%%%%%%%%%%%
%% The bootstrapped switch, shared by l5_sw2 and l5_sw3.
%%
%% l5_sw2 shows one of them next to a bulk switched transmission gate;
%% l5_sw3 shows two of them driving a differential pair with cross coupled
%% dummies. It is the same drawing all three times, so it lives here and the
%% figures inherit identical geometry.
%%
%% How it works. In phase c-bar the capacitor is strapped between ground and
%% V_DD and charges to the supply, while the gate is held at ground and the
%% switch is off. In phase c the bottom plate follows A and the top plate is
%% handed to the gate, so V_GS = V_DD whatever A does. That is the point: the
%% on resistance no longer depends on the signal.
%%
%% \bootBlockDn is the same block drawn below its rail rather than above it.
%% l5_sw3 needs it: with both halves pointing the same way the lower network
%% sits in the middle of the figure and the wires to the dummies have nowhere
%% to run, which is what makes the original crowded.
%%
%% Every name is prefixed. Short generic names collide with LaTeX built-ins --
%% \th is thorn -- and the redefinition is fatal, not a warning.
%%%%%%%%%%%%%%%%%%%%%%%%%%%%%%%%%%%%%%%%%%%%%%%%%%%%%%%%%%%%%%%%%%%%%%

% A clock-phase switch, hinge first, spanning 1.0 to the right.
\newcommand{\bootSwH}[3]{%
  \draw[thick] (#1,#2) -- (#1+0.15,#2);
  \fill (#1+0.15,#2) circle (0.055);
  \draw[thick] (#1+0.15,#2) -- (#1+0.88,#2+0.48);
  \fill (#1+0.85,#2) circle (0.055);
  \draw[thick] (#1+0.85,#2) -- (#1+1.0,#2);
  \node[anchor=north] at (#1+0.5,#2-0.12) {\small #3};
}

% The same, mirrored about its own rail: blade down, label above.
\newcommand{\bootSwHd}[3]{%
  \draw[thick] (#1,#2) -- (#1+0.15,#2);
  \fill (#1+0.15,#2) circle (0.055);
  \draw[thick] (#1+0.15,#2) -- (#1+0.88,#2-0.48);
  \fill (#1+0.85,#2) circle (0.055);
  \draw[thick] (#1+0.85,#2) -- (#1+1.0,#2);
  \node[anchor=south] at (#1+0.5,#2+0.12) {\small #3};
}

% The same, running 1.0 downward from (#1,#2).
\newcommand{\bootSwD}[3]{%
  \draw[thick] (#1,#2) -- (#1,#2-0.15);
  \fill (#1,#2-0.15) circle (0.055);
  \draw[thick] (#1,#2-0.15) -- (#1+0.48,#2-0.88);
  \fill (#1,#2-0.85) circle (0.055);
  \draw[thick] (#1,#2-0.85) -- (#1,#2-1.0);
  \node[anchor=east] at (#1-0.12,#2-0.5) {\small #3};
}

% The same, running 1.0 upward from (#1,#2).
\newcommand{\bootSwU}[3]{%
  \draw[thick] (#1,#2) -- (#1,#2+0.15);
  \fill (#1,#2+0.15) circle (0.055);
  \draw[thick] (#1,#2+0.15) -- (#1+0.48,#2+0.88);
  \fill (#1,#2+0.85) circle (0.055);
  \draw[thick] (#1,#2+0.85) -- (#1,#2+1.0);
  \node[anchor=east] at (#1-0.12,#2+0.5) {\small #3};
}

% Grounds, drawn by hand. circuitikz's `ground` node shape renders as nothing
% at all in the TeX Live this repository builds with -- the wire simply ends --
% so every figure here draws the bars itself, the way ckt_lib's \vground does.
% Bars to the left of the terminal, for a wire arriving from the right.
\newcommand{\bootGndL}[2]{%
  \draw[thick] (#1,#2-0.20) -- (#1,#2+0.20);
  \draw[thick] (#1-0.10,#2-0.13) -- (#1-0.10,#2+0.13);
  \draw[thick] (#1-0.20,#2-0.06) -- (#1-0.20,#2+0.06);
}
% Bars to the right, for a wire arriving from the left.
\newcommand{\bootGndR}[2]{%
  \draw[thick] (#1,#2-0.20) -- (#1,#2+0.20);
  \draw[thick] (#1+0.10,#2-0.13) -- (#1+0.10,#2+0.13);
  \draw[thick] (#1+0.20,#2-0.06) -- (#1+0.20,#2+0.06);
}
% Bars below, for a wire arriving from above.
\newcommand{\bootGndD}[2]{%
  \draw[thick] (#1-0.20,#2) -- (#1+0.20,#2);
  \draw[thick] (#1-0.13,#2-0.10) -- (#1+0.13,#2-0.10);
  \draw[thick] (#1-0.06,#2-0.20) -- (#1+0.06,#2-0.20);
}

% A labelled terminal: coordinate, anchor, label.
\newcommand{\bootPort}[3]{%
  \draw[thick] (#1) circle (0.09);
  \node[anchor=#2] at (#1) {#3};
}

% The whole bootstrapped switch, drawn above the input rail at y = 0.
% #1 is the name of the pass device, so a figure can place more than one.
% The input rail runs from x = -3.2, the output rail to x = 6.8; the caller
% adds the terminals.
\newcommand{\bootBlock}[1]{%
  % ---- Bottom plate: ground when off, the input when on ----
  \bootGndL{-1.6}{1.8}
  \draw[thick] (-1.6,1.8) -- (-1.4,1.8);
  \bootSwH{-1.4}{1.8}{$\overline{c}$}
  \draw[thick] (-0.4,1.8) -- (0.0,1.8);
  \fill (-0.2,1.8) circle (0.07);
  \bootSwD{-0.2}{1.6}{$c$}
  \draw[thick] (-0.2,0.6) -- (-0.2,0);
  \fill (-0.2,0) circle (0.07);
  % ---- The bootstrap capacitor ----
  \draw[thick] (0.0,1.8) to [C] (1.4,1.8);
  \draw[thick] (1.4,1.8) -- (2.4,1.8);
  % ---- Top plate: V_DD when off, the gate when on ----
  \fill (2.0,1.8) circle (0.07);
  \bootSwU{2.0}{2.0}{$\overline{c}$}
  \draw[thick, ->] (2.0,3.0) -- (2.0,3.45) node[anchor=south] {$V_{DD}$};
  \bootSwH{2.4}{1.8}{$c$}
  \draw[thick] (3.4,1.8) -- (4.0,1.8);
  % ---- Gate: grounded when off ----
  \fill (3.7,1.8) circle (0.07);
  \bootSwH{4.0}{1.8}{$\overline{c}$}
  \draw[thick] (5.0,1.8) -- (5.2,1.8);
  \bootGndR{5.2}{1.8}
  % ---- The pass device ----
  \draw[thick] (2.9,0) to [Tnmos, n=#1] (4.5,0);
  \draw[thick] (#1.gate) -- (3.7,1.8);
  \draw[thick] (-3.2,0) -- (2.9,0);
  \draw[thick] (4.5,0) -- (6.8,0);
}

% The same block drawn below its rail.
\newcommand{\bootBlockDn}[1]{%
  \bootGndL{-1.6}{-1.8}
  \draw[thick] (-1.6,-1.8) -- (-1.4,-1.8);
  \bootSwHd{-1.4}{-1.8}{$\overline{c}$}
  \draw[thick] (-0.4,-1.8) -- (0.0,-1.8);
  \fill (-0.2,-1.8) circle (0.07);
  \bootSwU{-0.2}{-1.6}{$c$}
  \draw[thick] (-0.2,-0.6) -- (-0.2,0);
  \fill (-0.2,0) circle (0.07);
  \draw[thick] (0.0,-1.8) to [C] (1.4,-1.8);
  \draw[thick] (1.4,-1.8) -- (2.4,-1.8);
  \fill (2.0,-1.8) circle (0.07);
  \bootSwD{2.0}{-2.0}{$\overline{c}$}
  \draw[thick, ->] (2.0,-3.0) -- (2.0,-3.45) node[anchor=north] {$V_{DD}$};
  \bootSwHd{2.4}{-1.8}{$c$}
  \draw[thick] (3.4,-1.8) -- (4.0,-1.8);
  \fill (3.7,-1.8) circle (0.07);
  \bootSwHd{4.0}{-1.8}{$\overline{c}$}
  \draw[thick] (5.0,-1.8) -- (5.2,-1.8);
  \bootGndR{5.2}{-1.8}
  \draw[thick] (2.9,0) to [Tnmos, n=#1, mirror] (4.5,0);
  \draw[thick] (#1.gate) -- (3.7,-1.8);
  \draw[thick] (-3.2,0) -- (2.9,0);
  \draw[thick] (4.5,0) -- (6.8,0);
}
